\documentclass[11pt]{article}
\usepackage[margin=1in]{geometry}
\usepackage{amsmath,amssymb,amsthm,booktabs,hyperref}
\newtheorem{theorem}{Theorem}
\newtheorem{lemma}[theorem]{Lemma}
\newcommand{\F}{\mathbb F}
\title{Local gcds of Hankel determinants in Shallit's four-letter word}
\author{Blake Stone's mathematics project\\\small Developed with OpenAI Codex}
\date{November 26, 2026}
\begin{document}\maketitle
\begin{abstract}
In the fixed word of Shallit's substitution the shifted Hankel determinants of any order have gcd one in every short interval of starting positions. At every fixed starting position the gcd over a short interval of orders is also one. Both estimates are uniform in the order. The proof uses pairs of equal left contexts whose next symbols differ by one.
\end{abstract}
Let $w_0w_1\ldots$ be the fixed point starting in $1$ of
$\phi(1)=12$, $\phi(2)=23$, $\phi(3)=14$, $\phi(4)=32$. Put
\[
 H_r(s)=\det(w_{s+i+j})_{0\le i,j<r},\qquad H_0(s)=1.
\]
Shallit asks whether $H_r(s)\ne0$ for every $r\ge1,s\ge0$ \cite{shallit,slides}. A previous repository calculation verifies all shifts up to order 4096. We prove a uniform arithmetic result for unbounded orders.

\begin{theorem}\label{main}
For all $s\ge0$ and $r\ge1$,
\begin{align*}
 \gcd\{H_r(s+i):0\le i<\lceil13r/2\rceil\}&=1,\\
 \gcd\{H_k(s):r\le k\le\lfloor15r/2\rfloor\}&=1 .
\end{align*}
In particular, for each prime $p$, each indicated interval contains a determinant nonzero modulo $p$.
\end{theorem}

\section{An obstruction to local linear recurrences}
Say that $D(d)$ is admissible if, for every $u\ge0$, the $D(d)$ endpoints beginning at $u+d$ include two having identical preceding $d$ letters and next letters differing by exactly one. Endpoints refer to the positions of the next letters. Over every field a fixed monic recurrence of degree $d$ cannot hold at all these endpoints. Indeed, the recurrence predicts the same next value from the same context, whereas the two actual values differ by a unit.

\begin{lemma}\label{bound}
There are admissible values $D(0)=4$, $D(1)=8$, and $D(d)\le 15d/2$ for every $d\ge2$.
\end{lemma}
\begin{proof}
The following two finite tables check contexts of length $h$ and sets of $b$ consecutive endpoints following them. An entry in row I means every length-$(h+b)$ factor has two equal $h$-contexts with next symbols differing by one. In row II one next letter must be $2$ and the other must be one of $1,3,4$.
\[
\begin{array}{c|rrrrrrrrr}\toprule
h&0&1&2&3&4&5&6&7&8\\\midrule
{\rm I}: b&4&8&13&18&20&36&36&40&-\\
{\rm II}: b&-&-&-&-&20&30&32&40&40\\\bottomrule
\end{array}
\]
The adjacent pairs in $w$ are
$12,14,21,22,23,31,32,41,42$. Each occurs in $\phi^5(1)$. The set contains the four substitution images and is closed under taking the middle pair of $\phi(ab)$ for each listed $ab$. For $t=h+b$ take $E$ to be the least power of two at least $t$ and write $E=2^j$. For each of the nine pairs inspect the $E$ length-$t$ factors of $\phi^j(ab)$ starting at offsets $0,\ldots,E-1$. This covers every length-$t$ factor by alignment in $w=\phi^j(w)$. For each such factor inspect all endpoint pairs among its last $b$ positions. The accompanying short program performs exactly these checks, with at most 329 distinct factors for a table entry.

For $d\ge8$ choose a power of two $B=2^k$ with $4B\le d<8B$, and put $h=\lceil d/B\rceil$. Then $k\ge1$. The first letter of $\phi^k(2)$ is $2$, while the first of each other image is $1$ or $3$. Put $J=\lceil(u+d)/B\rceil\ge h$. Use row II in the base word at the $b$ endpoints beginning at $J$, retaining the $h$ positions immediately before them for their contexts. Inflating gives endpoints at multiples of $B$ with equal preceding $hB\ge d$ letters and next letters differing by one. Their positions range from $u+d$ to at most $u+d+bB-1$. Thus $D(d)=bB$ is admissible. For $h=4$ this is $20B\le15d/2$. For $h=5,6,7,8$, respectively, the ratios $b/(h-1)$ are $30/4,32/5,40/6,40/7$. Each is at most $15/2$, and $d>(h-1)B$. Row I handles $d<8$.
\end{proof}

\section{Orders at one shift}
Fix a prime $p$ and shift $s$, and take all linear algebra in $\F_p$. Let the linear functional $L$ on $\F_p[X]$ be given by $L(X^i)=w_{s+i}$. If $H_d(s)\ne0$, there is a unique monic $q$ of degree $d$ satisfying
$L(qX^i)=0$ for $0\le i<d$, by solving a system with Hankel coefficient matrix. For $d=0$ take $q=1$. By Lemma~\ref{bound} the first $j$ with $L(qX^j)\ne0$ exists and $d\le j<D(d)$.

For $d<k\le j$ the coefficient vector of $q$ padded with zeros lies in the kernel of the order-$k$ Hankel matrix. At $k=j+1$ use the monic triangular basis
\[
 1,X,\ldots,X^{d-1},q,Xq,\ldots,X^{j-d}q
\]
for the bilinear form $(f,g)\mapsto L(fg)$. Its two displayed groups are orthogonal. The first block has determinant $H_d(s)$. In the second block the entries with indices $a,b$ are $L(q^2X^{a+b})$. They vanish for $a+b<j-d$ and equal $L(qX^j)\ne0$ on $a+b=j-d$. Reversing its columns gives a triangular nonsingular matrix. Thus the next nonsingular order after $d$ is $j+1\le D(d)$.

Starting from zero this produces arbitrarily large nonsingular orders. For a prescribed $r\ge1$, take the first such order $k\ge r$ and its predecessor $d<r$. By Lemma~\ref{bound}, $D(d)\le\lfloor15r/2\rfloor$, including $d=0,1$. This proves the fixed-shift assertion modulo $p$.

\section{Shifts at one order}
Suppose $K=\lceil13r/2\rceil$ consecutive determinants of order $r$ are zero modulo $p$. Descend through the smaller orders. Whenever the $K$ consecutive determinants $H_{d+1}(a),\ldots,H_{d+1}(a+K-1)$ vanish, the Desnanot--Jacobi identity (or $H_0=1$ for $d=0$) gives
\[
 H_d(x)H_d(x+2)=H_d(x+1)^2,\qquad a\le x<a+K.
\]
If the $K$ middle entries $H_d(a+1),\ldots,H_d(a+K)$ all vanish, continue from this interval of length $K$. Otherwise any nonzero middle entry forces both its neighbors nonzero, and repeated use of the identity makes all $K+2$ entries nonzero. The descent reaches this second alternative by $d=0$.

For $d\ge1$, fit the unique monic recurrence of degree $d$ at each start $x=a,\ldots,a+K-1$, using the first $2d$ letters there. Since $H_{d+1}(x)=0$ and $H_d(x)\ne0$, the recurrence also holds at the next endpoint $x+2d$ by the Schur complement. The recurrence at $x$ thus solves the $d$ fitting equations at $x+1$, and uniqueness follows from $H_d(x+1)\ne0$. All the fitted recurrences coincide. This yields a single recurrence at $K+d$ consecutive endpoints beginning at $a+d$. For $d=0$ the same conclusion is $K$ consecutive zero letters. In either case $K+d\ge D(d)$: check $d=0,1$ directly and use $d<r$ with Lemma~\ref{bound} for $d\ge2$. This contradicts the identical-context obstruction. We have shown that no prime divides all the entries of either interval, proving Theorem~\ref{main}. The nonvanishing of each individual integer determinant remains open.

\begin{thebibliography}{9}
\bibitem{shallit} J. Shallit, On patterns and pattern avoidance, in \emph{Mini-Workshop: Combinatorics on Words}, Oberwolfach Reports 7 (2010), 2234, Problems 44--45. \href{https://ems.press/content/serial-article-files/46296}{Workshop report}.
\bibitem{slides} J. Shallit, \emph{Open Problems in Automata Theory: An Idiosyncratic View}, BCTCS 2014, slides 33--35. \href{https://cs.uwaterloo.ca/~shallit/Talks/bc4.pdf}{Lecture slides}.
\bibitem{lunnon} F. Lunnon, The Number-Wall Algorithm: An LFSR Cookbook, J. Integer Seq. 4 (2001), Article 01.1.1. \href{https://cs.uwaterloo.ca/journals/JIS/VOL4/LUNNON/numbwall10.html}{Journal}.
\end{thebibliography}
\end{document}
